\section{Complete-contrast inference for the existing fitted scores}
\label{sec:joint-contrast}

Keep the fitted scores and sample allocation from the preceding refinement.
The source center is $\hat T_S=\hat d+K^{-1}\sum_j(\bar R_j^1-\bar R_j^0)$.
The change here concerns its noise and nuisance certificates, not its weights.
Two source arms need not use the same valid set.

\paragraph{A joint patient-noise bound.}
Group $g$ contains $n_g$ independent evaluation patients, group mean $\bar U_g$,
and unbiased mean variance $\hat v_g=S_g^2/n_g$. All centers $c_g$, weights
$w_g$, and support bounds $[l_g,u_g]$ are fixed by training. Let $N=\sum_gn_g$,
\begin{align*}
 S&=\sum_gw_g(\bar U_g-c_g),\\
 A_2&=\sum_gw_g^2\left\{\frac{n_g-1}{n_g}\hat v_g+
                         \frac{(\bar U_g-c_g)^2}{n_g}\right\},\\
 \hat V&=\frac{NA_2-S^2}{N-1},\qquad
 B=\max_{g,z\in\{l_g,u_g\}}\frac{w_g(z-c_g)}{n_g}
   -\min_{g,z\in\{l_g,u_g\}}\frac{w_g(z-c_g)}{n_g}.
\end{align*}
Theorem 11 of \citet{maurer2009}, applied to each patient variable
$Nw_g(U_{gi}-c_g)/n_g$, gives the two-sided radius
\[
 r_N=\sqrt{2\hat V\log(4/\alpha_N)}+
          \frac{7NB\log(4/\alpha_N)}{3(N-1)}.
\]
Its independence assumption permits nonidentical means. Consequently the
empirical variance above includes between-group centered means; it is not
silently replaced by only within-group variances.

For the joint TATE direction, one group is the target prediction contrast,
with weight one, and source group $j$ contains $R_j^1-R_j^0$, with weight
$1/K$. Training residual means supply the source centers; training prediction
means supply the target center. These offsets change the certificate, not
$\hat T_S$. They are not chosen using evaluation observations.

For the original patient residual scores, $R_i^1R_i^0=0$, but
$\widehat{\Cov}(\bar R^1,\bar R^0)=-\bar R^1\bar R^0/(n-1)$ need not vanish.
The contrast mean variance is therefore
$\hat v_1+\hat v_0-2\widehat{\Cov}(\bar R^1,\bar R^0)$.
Our ablation compares this joint bound with three separately protected means,
using the same total noise allowance.

\paragraph{A joint nuisance direction.}
Use the preceding design coefficient bounds $C_{V_a,x}$. A single application
of \eqref{eq:v4-scalar-radius} to the stacked coefficients
$(c_{V_1,1},\ldots,c_{V_1,L},-c_{V_0,1},\ldots,-c_{V_0,L})$ bounds
$|\bar\beta_{V_1}^1-\bar\beta_{V_0}^0|$. Conditional on the designs, these
are one linear combination of independent target training outcomes.
This avoids adding two separate standard-deviation bounds. It still assumes
outcome-independent weighting fits and protects a fixed population-valid
pair of subsets, not an outcome-selected pair.

\paragraph{Transport dispersion remains a distinct choice.}
The arm certificate is
\[
 r_D^{\rm arm}=\sqrt{(K-g_1)U_1/g_1}+\sqrt{(K-g_0)U_0/g_0}.
\]
Alternatively, if $g_\cap=g_1+g_0-K>0$, at least $g_\cap$ sites are valid in
both arms. Apply dispersion calibration to $R^1-R^0$ and its proper variance,
obtaining $U_\tau$. Its certificate is
$\sqrt{(K-g_\cap)U_\tau/g_\cap}$, with nuisance coefficients recomputed for
$g_\cap$ rather than $g_a$. It can exploit opposite terms that cancel in
TATE, but the smaller validity guarantee can increase nuisance protection.
The implementation rejects this mode when $g_\cap\le0$.
These are separately calibrated options; picking their shorter realized
interval without simultaneous calibration would not be justified.

\paragraph{Reporting and error allocation.}
The source band is $\hat T_S\pm(r_N+r_B+r_D)$. Use $.025$ for noise,
$.005$ for nuisance, $.01$ for dispersion, and $.01$ for a target-only
population interval. Report their intersection if nonempty, otherwise the
shorter band. The union bound gives at least .95 coverage under the stated
contracts. Both the source band and the final band are saved: small-$K$
results can otherwise conceal reliance on the target band.

A further comparator estimates an upper bound for the true target CATE range
from simultaneous full-target cell-outcome intervals. A range allowance .005,
plus outcome and composition allowances .0225 each, gives another .95
stratified target-only interval. The random range bound is used only on its
simultaneous certificate event; the composition concentration inequality itself
is for the fixed true CATE function. This is a valid refinement of the earlier
worst-case range-two comparator.
